\section*{17. 数组最大子段的个数（O(n)）}

设$A[1..n]$是长度为$n$的数组，存放可重复的整数。最大子段指其和不小于任何子段和的子段。要求设计$O(n)$算法，求最大子段的个数。



\begin{itemize}
    \item \textbf{状态定义}：$currSum$当前和，$currCount$当前个数，$globalSum$全局最大，$globalCount$全局个数
    \item \textbf{转移规则}：若$A[i] > currSum + A[i]$则新子段；若$A[i] < currSum + A[i]$则延续；若相等则个数相加
    \item \textbf{全局更新}：若$currSum > globalSum$则更新；若相等则累加个数
    \item \textbf{时间复杂度}：$O(n)$
\end{itemize}
\begin{lstlisting}[language=C++, style=cppstyle]
int maxSubarrayCount(vector<int>& A) {
    int n = A.size();
    if (n == 0) return 0;

    long long currSum = A[0];
    long long currCount = 1;
    long long globalSum = A[0];
    long long globalCount = 1;

    for (int i = 1; i < n; i++) {
        long long sumExtend = currSum + A[i];
        long long sumNew = A[i];

        if (sumNew > sumExtend) {
            // 新子段更优
            currSum = sumNew;
            currCount = 1;
        } else if (sumNew < sumExtend) {
            // 延续更优
            currSum = sumExtend;
        } else {
            // 两者相等，个数相加
            currSum = sumExtend;
            currCount += currCount;
        }

        // 更新全局最大值
        if (currSum > globalSum) {
            globalSum = currSum;
            globalCount = currCount;
        } else if (currSum == globalSum) {
            globalCount += currCount;
        }
    }

    return globalCount;
}
\end{lstlisting}


\begin{enumerate}
    \item 维护"当前和"、"当前个数"、"全局最大"、"全局个数"
    \item 比较"重新开始"和"延续前面"哪个更大
    \item 相等则累加个数
\end{enumerate}

